\section{Conclusion}
\label{sec:conclusion}

The displacement formulation distinguishes the magnitude of individual updates,
their local directional interaction and their cumulative accumulation. For
plain SGD, the exact decomposition retains all temporal lags, while consecutive
cosine provides a tractable local observation. The curvature and sampling
relations clarify how that observation can change without determining either
cumulative displacement or loss improvement. The full-batch and mini-batch
experiments illustrate these distinctions: gradient reversal can coexist with
objective progress, and weaker reversal need not improve displacement retention
or the attained loss.

Under Adam, raw-gradient diagnostics remain informative about the loss
landscape traversed by adaptive steps. Controlled full-batch experiments
connect loss spikes with a predominantly destructive pre-rise gradient
regime and sharp norm amplification. The micro-plots resolve projected
gradient changes before the loss maximum; the event population quantifies
the variation in timing and recovery across models and rates. A joint
descriptor of pre-rise reversal and rise norm growth is associated with
relative peak magnitude, without establishing a causal feedback mechanism
or a universal instability threshold.

The practical contribution is a framework for interpreting these measurements
together and using them to formulate controlled interventions. Developing a
scheduler or clipping rule requires demonstrating that acting on early
diagnostics improves subsequent loss at matched budgets, beyond decisions
based on recent loss and known training settings. Prospective validation,
controlled interventions and replication across independent runs remain
future work.
